\documentclass[11pt,a4paper]{article}

\usepackage[margin=2.2cm]{geometry}
\usepackage{amsmath,amssymb,amsthm}
\usepackage{braket}
\usepackage{physics}
\usepackage{booktabs}
\usepackage{enumitem}
\usepackage{xcolor}
\usepackage{tikz}
\usetikzlibrary{quantikz}

\newcommand{\N}{\mathbb{N}}
\newcommand{\R}{\mathbb{R}}
\newcommand{\WMC}{\operatorname{WMC}}

\title{Quantum Computing Mock Examination\\Complete Solutions}
\author{Summer Semester 2026}
\date{}

\begin{document}
\maketitle

\section*{Question 1: Reversible Computing and Toffoli Gates}

\subsection*{(a) Reversible circuit for $f(a,b,c)=(a\land b)\lor c$}
Use five wires $a,b,c,x,y$, where initially $x=y=0$. First compute
$x=a\land b$, and then use
\[
A\lor B=A\oplus B\oplus(A\land B).
\]
Thus the following sequence computes the desired output:
\begin{enumerate}[label=\arabic*.]
  \item $\operatorname{Toffoli}(a,b;x)$, so that $x\leftarrow a\land b$.
  \item $\operatorname{CNOT}(c;y)$, so that $y\leftarrow c$.
  \item $\operatorname{CNOT}(x;y)$, so that $y\leftarrow c\oplus x$.
  \item $\operatorname{Toffoli}(x,c;y)$, so that
  \[
  y\leftarrow c\oplus x\oplus(cx)=c\lor x=(a\land b)\lor c.
  \]
\end{enumerate}
A circuit diagram is
\[
\begin{quantikz}
\lstick{$a$}     & \ctrl{3} & \qw      & \qw      & \qw      & \qw \\
\lstick{$b$}     & \ctrl{2} & \qw      & \qw      & \qw      & \qw \\
\lstick{$c$}     & \qw      & \ctrl{2} & \qw      & \ctrl{2} & \qw \\
\lstick{$x=0$}   & \targ{}  & \qw      & \ctrl{1} & \ctrl{1} & \qw \\
\lstick{$y=0$}   & \qw      & \targ{}  & \targ{}  & \targ{}  & \rstick{$f(a,b,c)$}
\end{quantikz}
\]
All input bits are preserved, so the overall transformation is reversible.

\subsection*{(b) Matrix of the Toffoli gate}
The action is
\[
\ket{a,b,c}\longmapsto\ket{a,b,c\oplus(a\land b)}.
\]
Only $\ket{110}$ and $\ket{111}$ are exchanged. Hence, in the stated basis order,
\[
P_{\mathrm{Toff}}=
\begin{pmatrix}
1&0&0&0&0&0&0&0\\
0&1&0&0&0&0&0&0\\
0&0&1&0&0&0&0&0\\
0&0&0&1&0&0&0&0\\
0&0&0&0&1&0&0&0\\
0&0&0&0&0&1&0&0\\
0&0&0&0&0&0&0&1\\
0&0&0&0&0&0&1&0
\end{pmatrix}.
\]
Since this matrix is real and symmetric,
\[
P_{\mathrm{Toff}}^T=P_{\mathrm{Toff}}.
\]
Moreover, applying Toffoli twice restores every basis state, so
\[
P_{\mathrm{Toff}}^2=I_8.
\]
Therefore
\[
P_{\mathrm{Toff}}^T P_{\mathrm{Toff}}=I_8.
\]

\subsection*{(c) Uncomputing the ancilla $x$}
After the output $y$ has been computed, apply the Toffoli gate with controls $a,b$
and target $x$ once more:
\[
x\leftarrow x\oplus(a\land b)
=(a\land b)\oplus(a\land b)=0.
\]
This does not affect $a,b,c$, or $y$.

In general, if a computation is
\[
U=U_kU_{k-1}\cdots U_1,
\]
then its inverse is
\[
U^{-1}=U_1^{-1}U_2^{-1}\cdots U_k^{-1}.
\]
Thus the inverse gates must be applied in reverse temporal order. This is necessary
because matrix multiplication is generally noncommutative.

\subsection*{(d) No-erasure theorem}
A reversible gate on $n$ bits is a bijection on $\{0,1\}^n$, equivalently a
$2^n\times 2^n$ permutation matrix.

Assume for contradiction that a reversible gate $G$ sets the last output bit to $0$
for every input. Then all outputs belong to
\[
S=\{z\in\{0,1\}^n:z_n=0\},
\]
which has only $2^{n-1}$ elements. However, bijectivity requires the $2^n$ different
inputs to have $2^n$ different outputs. This is impossible because
$2^{n-1}<2^n$. Hence no reversible gate can erase a bit unconditionally.

\newpage
\section*{Question 2: Markov Chains}

Let
\[
P=
\begin{pmatrix}
0&\tfrac12&0&\tfrac12\\
\tfrac14&\tfrac12&\tfrac14&0\\
0&0&\tfrac12&\tfrac12\\
\tfrac12&0&\tfrac14&\tfrac14
\end{pmatrix}.
\]

\subsection*{(a) Irreducibility}
The directed graph of positive-probability transitions contains
\[
0\to1,\quad 0\to3,\quad 1\to0,\quad 1\to2,
\quad 2\to3,\quad 3\to0,\quad 3\to2.
\]
For example,
\[
0\to1\to2,\qquad 2\to3\to0,\qquad 3\to0\to1.
\]
Thus every state can reach every other state. Hence the chain is irreducible.

\subsection*{(b) Aperiodicity}
Since
\[
P_{11}=\frac12>0,
\]
state $1$ can return to itself in one step. Therefore
\[
d(1)=\gcd\{t\ge1:P^{(t)}_{11}>0\}=1.
\]
In an irreducible chain all states have the same period, so every state has period $1$.
The chain is therefore aperiodic.

\subsection*{(c) Distribution after two steps}
With row-vector convention,
\[
p(0)=(1,0,0,0),
\]
so
\[
p(1)=p(0)P=(0,\tfrac12,0,\tfrac12).
\]
Then
\begin{align*}
p_0(2)&=\frac12\frac14+\frac12\frac12=\frac38,\\
p_1(2)&=\frac12\frac12+\frac12\cdot0=\frac14,\\
p_2(2)&=\frac12\frac14+\frac12\frac14=\frac14,\\
p_3(2)&=\frac12\cdot0+\frac12\frac14=\frac18.
\end{align*}
Thus
\[
\boxed{p(2)=\left(\frac38,\frac14,\frac14,\frac18\right)}.
\]

\subsection*{(d) Stationary distribution}
Let $\pi=(\pi_0,\pi_1,\pi_2,\pi_3)$. The equation $\pi P=\pi$ gives
\begin{align*}
\pi_0&=\frac14\pi_1+\frac12\pi_3,\\
\pi_1&=\frac12\pi_0+\frac12\pi_1,\\
\pi_2&=\frac14\pi_1+\frac12\pi_2+\frac14\pi_3,\\
\pi_3&=\frac12\pi_0+\frac12\pi_2+\frac14\pi_3.
\end{align*}
From the second equation, $\pi_1=\pi_0$. The first equation then yields
\[
\pi_3=\frac32\pi_0.
\]
The third equation gives
\[
\pi_2=\frac12(\pi_1+\pi_3)=\frac54\pi_0.
\]
Using normalization,
\[
\pi_0+\pi_1+\pi_2+\pi_3
=\pi_0\left(1+1+\frac54+\frac32\right)
=\frac{19}{4}\pi_0=1.
\]
Therefore
\[
\boxed{\pi=\left(\frac4{19},\frac4{19},\frac5{19},\frac6{19}\right)}.
\]

\newpage
\section*{Question 3: Quantum States and Circuit Analysis}

\subsection*{(a) The state $\ket{110}$}
\[
\ket{110}=\ket1\otimes\ket1\otimes\ket0
=
\begin{pmatrix}0\\1\end{pmatrix}
\otimes
\begin{pmatrix}0\\1\end{pmatrix}
\otimes
\begin{pmatrix}1\\0\end{pmatrix}.
\]
The only nonzero component is at binary index $110_2=6$, so
\[
\boxed{
\ket{110}=
\begin{pmatrix}
0\\0\\0\\0\\0\\0\\1\\0
\end{pmatrix}}
\]
in the basis $\ket{000},\ldots,\ket{111}$.

\subsection*{(b) Computing $(H\otimes H)\ket{10}$}
First,
\[
H\ket1=\frac1{\sqrt2}\begin{pmatrix}1\\-1\end{pmatrix}=\ket-,
\qquad
H\ket0=\frac1{\sqrt2}\begin{pmatrix}1\\1\end{pmatrix}=\ket+.
\]
Therefore
\begin{align*}
(H\otimes H)\ket{10}
&=(H\ket1)\otimes(H\ket0)\\
&=\frac12
\begin{pmatrix}1\\-1\end{pmatrix}
\otimes
\begin{pmatrix}1\\1\end{pmatrix}\\
&=\boxed{\frac12
\begin{pmatrix}1\\1\\-1\\-1\end{pmatrix}}.
\end{align*}
Directly,
\[
H\otimes H=\frac12
\begin{pmatrix}
1&1&1&1\\
1&-1&1&-1\\
1&1&-1&-1\\
1&-1&-1&1
\end{pmatrix},
\qquad
\ket{10}=\begin{pmatrix}0\\0\\1\\0\end{pmatrix}.
\]
Multiplication selects the third column, again giving
$\frac12(1,1,-1,-1)^T$.

\subsection*{(c) Reversing the direction of CNOT with Hadamards}
Let
\[
A=H\otimes H=\frac12
\begin{pmatrix}
1&1&1&1\\
1&-1&1&-1\\
1&1&-1&-1\\
1&-1&-1&1
\end{pmatrix}
\]
and
\[
C_{1\to2}=
\begin{pmatrix}
1&0&0&0\\
0&1&0&0\\
0&0&0&1\\
0&0&1&0
\end{pmatrix}.
\]
A direct multiplication gives
\[
A C_{1\to2} A
=
\begin{pmatrix}
1&0&0&0\\
0&0&0&1\\
0&0&1&0\\
0&1&0&0
\end{pmatrix}
=C_{2\to1}.
\]
Hence
\[
\boxed{(H\otimes H)C_{1\to2}(H\otimes H)=C_{2\to1}}.
\]

\subsection*{(d) Why quantum gates must be unitary}
For a normalized state $\ket\psi$,
\[
\braket{\psi|\psi}=1.
\]
After applying $U$, the state is $\ket{\psi'}=U\ket\psi$. Its norm is
\[
\braket{\psi'|\psi'}
=\bra\psi U^\dagger U\ket\psi.
\]
To preserve total probability for every input state, this must equal
$\braket{\psi|\psi}$ for all $\ket\psi$. Therefore
\[
U^\dagger U=I.
\]
Thus $U$ must be unitary. Unitarity also implies $U^{-1}=U^\dagger$, so the evolution is reversible.

\newpage
\section*{Question 4: Grover's Search Algorithm}

Let $N=8$ and let $\ket{x^*}$ be the unique marked state. Define
\[
\ket{r}=\frac1{\sqrt{N-1}}\sum_{x\ne x^*}\ket{x}.
\]
Then $\braket{x^*|r}=0$.

\subsection*{(a) Initial decomposition}
The uniform state is
\[
\ket{s}=\frac1{\sqrt N}\ket{x^*}
+\sqrt{\frac{N-1}{N}}\ket r.
\]
For $N=8$,
\[
\boxed{\ket{s}=\frac1{\sqrt8}\ket{x^*}+\sqrt{\frac78}\ket r}.
\]
The initial marked-state amplitude is $1/\sqrt8$.

\subsection*{(b) Phase kickback}
Prepare an ancilla in
\[
\ket- =\frac{\ket0-\ket1}{\sqrt2}
\]
and use the standard oracle
\[
U_f\ket{x}\ket y=\ket{x}\ket{y\oplus f(x)}.
\]
Then
\begin{align*}
U_f\ket{x}\ket-
&=\ket{x}\frac{\ket{f(x)}-\ket{1\oplus f(x)}}{\sqrt2}\\
&=(-1)^{f(x)}\ket{x}\ket-.
\end{align*}
The ancilla returns to the same state $\ket-$; only a phase is transferred to the input register.

\subsection*{(c) Diffusion operator and one Grover iteration}
The diffusion operator is
\[
D=2\ket{s}\bra{s}-I.
\]
Certainly,
\[
D\ket{s}=\ket{s}.
\]
However, with the definition above, $\ket r$ is orthogonal to $\ket{x^*}$ but
\emph{not} orthogonal to $\ket s$, because
\[
\braket{s|r}=\sqrt{\frac{N-1}{N}}\ne0.
\]
Therefore the statement $D\ket r=-\ket r$ is false. The correct action in the basis
$\{\ket{x^*},\ket r\}$ is
\[
D=
\begin{pmatrix}
\frac{2}{N}-1 & \frac{2\sqrt{N-1}}{N}\\[3pt]
\frac{2\sqrt{N-1}}{N} & \frac{N-2}{N}
\end{pmatrix}.
\]

For
\[
\ket\psi=\alpha\ket{x^*}+\beta\ket r,
\]
the oracle produces
\[
O_f\ket\psi=-\alpha\ket{x^*}+\beta\ket r.
\]
Applying $D$, the new marked-state amplitude is
\[
\boxed{
\alpha'=\frac{N-2}{N}\alpha+\frac{2\sqrt{N-1}}{N}\beta}.
\]
For $N=8$,
\[
\boxed{\alpha'=\frac34\alpha+\frac{\sqrt7}{4}\beta}.
\]

Equivalently, writing $\sin\theta=1/\sqrt N$, one Grover iteration rotates the state
by $2\theta$ in the two-dimensional space spanned by $\ket{x^*}$ and $\ket r$.

\subsection*{(d) Optimal iteration count and success probability}
Let
\[
\theta=\arcsin\left(\frac1{\sqrt8}\right)\approx0.361367.
\]
After $k$ iterations, the marked-state amplitude is
\[
\sin((2k+1)\theta),
\]
so the success probability is
\[
P_k=\sin^2((2k+1)\theta).
\]
The real-valued optimum is approximately
\[
k_*\approx\frac{\pi}{4\theta}-\frac12\approx1.673,
\]
so the best integer is $k=2$. Then
\[
P_2=\sin^2(5\theta)\approx0.9453.
\]
Thus
\[
\boxed{k=2,\qquad P_{\mathrm{success}}\approx94.5\%.}
\]

\newpage
\section*{Question 5: Quantum Circuit Simulation via Weighted \#SAT}

\subsection*{(a) Boolean encodings of reversible gates}
\begin{enumerate}[label=(\roman*)]
  \item NOT:
  \[
  F_{\mathrm{NOT}}(x,x')=(x'\leftrightarrow\neg x).
  \]
  \item CNOT:
  \[
  F_{\mathrm{CNOT}}=(x'\leftrightarrow x)
  \land(y'\leftrightarrow(y\oplus x)).
  \]
  \item Toffoli:
  \[
  F_{\mathrm{Toff}}=(x'\leftrightarrow x)
  \land(y'\leftrightarrow y)
  \land(z'\leftrightarrow(z\oplus(x\land y))).
  \]
\end{enumerate}

\subsection*{(b) Hadamard encoding}
Use
\[
F_H(x,x',h)=\bigl(h\leftrightarrow(x\land x')\bigr),
\]
with literal weights
\[
W(h)=-\frac1{\sqrt2},
\qquad
W(\neg h)=\frac1{\sqrt2}.
\]
For fixed values of $x$ and $x'$, the constraint fixes $h$ uniquely:
\[
\begin{array}{c c c c}
\toprule
x&x'&h&\WMC\\
\midrule
0&0&0&+1/\sqrt2\\
0&1&0&+1/\sqrt2\\
1&0&0&+1/\sqrt2\\
1&1&1&-1/\sqrt2\\
\bottomrule
\end{array}
\]
These are exactly the entries $H_{x',x}$.

The auxiliary variable is required to carry the sign information. A plain Boolean
constraint on $x,x'$ can indicate whether a transition is allowed, but without a
weighted auxiliary literal it cannot distinguish $+1/\sqrt2$ from $-1/\sqrt2$.

\subsection*{(c) Bell-state circuit}
The complete formula is
\begin{align*}
F_{\mathrm{circuit}}={}&(\neg x_0)\land(\neg y_0)\\
&\land\bigl(h\leftrightarrow(x_0\land x_1)\bigr)\\
&\land(y_1\leftrightarrow y_0)\\
&\land(x_2\leftrightarrow x_1)\\
&\land\bigl(y_2\leftrightarrow(y_1\oplus x_1)\bigr).
\end{align*}
The weights are
\[
W(h)=-\frac1{\sqrt2},\qquad W(\neg h)=\frac1{\sqrt2},
\]
and all other literals have weight $1$.

Because $x_0=y_0=0$, the Hadamard constraint forces $h=0$ for both possible values
of $x_1$. Also $y_1=0$, $x_2=x_1$, and $y_2=x_1$. Therefore the only output paths are
\[
(x_2,y_2)=(0,0)\quad\text{and}\quad(1,1),
\]
each with weight $1/\sqrt2$. Hence
\[
\alpha_{00}=\frac1{\sqrt2},
\qquad
\alpha_{11}=\frac1{\sqrt2},
\qquad
\alpha_{01}=\alpha_{10}=0.
\]
The measurement probabilities are
\[
\boxed{P(00)=P(11)=\frac12.}
\]

\newpage
\section*{Question 6: Quantum Complexity and Variational Algorithms}

\subsection*{(a) BQP and QMA}
\textbf{BQP} is the class of decision problems solvable by a uniform family of
polynomial-size quantum circuits with bounded error, conventionally at most $1/3$.
Its randomized classical analogue is BPP.

\textbf{QMA} is the class of decision problems for which every yes-instance has a
polynomial-size quantum witness that can be verified by a polynomial-time quantum
verifier with bounded completeness and soundness error. Its closest classical analogue
is NP, more precisely the probabilistic verifier class MA when comparing verifier structure.

The standard inclusions are
\[
P\subseteq BPP\subseteq BQP\subseteq PSPACE.
\]
The first inclusion holds because deterministic computation is a special case of
randomized computation. The second holds because a quantum computer can simulate
classical probabilistic computation. The final inclusion follows because polynomial-size
quantum computations can be simulated using polynomial classical space, although usually
with exponential time.

\subsection*{(b) Meaning of $\mathrm{QPSPACE}=\mathrm{PSPACE}$}
This equality means that allowing quantum computation does not enlarge the class of
problems solvable with polynomial space. It is a complexity-class statement: every
polynomial-space quantum computation can be simulated in polynomial classical space,
and classical polynomial-space computation is trivially simulable quantumly.

It does \emph{not} rule out polynomial-factor or constant-factor memory savings on
particular problems; it says there is no separation at the level of polynomial-space
versus super-polynomial-space computability.

\subsection*{(c) QAOA for Max-Cut}
For an edge $(i,j)$, define
\[
C_{ij}=\frac{I-Z_iZ_j}{2}.
\]
This has eigenvalue $1$ when the endpoint bits differ and $0$ when they agree. Thus
\[
H_C=\sum_{(i,j)\in E}\frac{I-Z_iZ_j}{2}.
\]
The mixing Hamiltonian is
\[
H_M=\sum_{i\in V}X_i.
\]
Starting from
\[
\ket{s}=\ket+^{\otimes n},
\]
the depth-$1$ ansatz is
\[
\boxed{
\ket{\psi(\gamma,\beta)}
=e^{-i\beta H_M}e^{-i\gamma H_C}\ket{s}.}
\]
The classical optimizer tunes $\gamma$ and $\beta$ to maximize
\[
\langle H_C\rangle_{\gamma,\beta}
=\bra{\psi(\gamma,\beta)}H_C\ket{\psi(\gamma,\beta)}.
\]
Equivalently, it may minimize $-\langle H_C\rangle$. Difficulties include a
nonconvex parameter landscape, sampling noise, hardware noise, and in some regimes
small gradients or barren plateaus.

\subsection*{(d) Deutsch--Jozsa algorithm}
The input is initialized as
\[
\ket0^{\otimes n}\ket1.
\]
The ancilla starts in $\ket1$ because the first Hadamard transforms it into
\[
H\ket1=\ket-.
\]
After applying $H^{\otimes(n+1)}$, the state is
\[
\frac1{\sqrt{2^n}}\sum_{x\in\{0,1\}^n}\ket{x}\ket-.
\]
The oracle gives phase kickback:
\[
U_f\ket{x}\ket-=(-1)^{f(x)}\ket{x}\ket-,
\]
so the state becomes
\[
\frac1{\sqrt{2^n}}\sum_x(-1)^{f(x)}\ket{x}\ket-.
\]
After the final $H^{\otimes n}$, the amplitude of $\ket{0^n}$ is
\[
\frac1{2^n}\sum_x(-1)^{f(x)}.
\]
If $f$ is constant, this amplitude is $+1$ or $-1$, hence the outcome $0^n$ occurs
with probability $1$. If $f$ is balanced, half the terms are $+1$ and half are $-1$,
so the amplitude is $0$. Therefore measuring $0^n$ means constant, while any other
outcome means balanced.

\end{document}
